\documentclass[12pt, a4paper]{article}

% Paquetes requeridos para el circuito
\usepackage[utf8]{inputenc}
\usepackage[T1]{fontenc}
\usepackage{circuitikz} 

\begin{document}

\begin{center}
    \begin{circuitikz}[american, scale=1.0, transform shape]
        
        % 1. Fuentes V1 e Io en serie (Rama izquierda)
        \draw (0,0) to[I, l=$I_o$] (0,2)
                    to[V, l=$V_1$] (0,4);
        
        % 2. Resistencias R1 y Ro
        \draw (0,2) to[R, l=$R_1$] (3,2);
        \draw (0,4) to[R, l=$R_o$] (3,4);
        
        % 3. Rail inferior (Conexión común del nodo 'b')
        \draw (0,0) -- (5,0);
        
        % 4. Resistencias R2 y R3 desde el extremo libre de R1
        \draw (3,2) to[R, l=$R_2$] (3,4);
        \draw (3,2) to[R, l=$R_3$] (3,0);
        
        % 5. Conexión desde la unión R2-R3 hacia el nodo 'a'
        \draw (3,2) -- (5,2);
        
        % 6. Fuente de corriente desde 'a' hacia la unión R2-Ro
        \draw (5,2) to[I] (5,4) -- (3,4);
        
        % 7. Resistencia R4 entre los nodos 'a' y 'b'
        \draw (5,2) to[R, l=$R_4$] (5,0);
        
        % 8. Terminales de salida 'a' y 'b'
        \draw (5,2) to[short, -o] (6,2) node[right] {a};
        \draw (5,0) to[short, -o] (6,0) node[right] {b};
        
    \end{circuitikz}
\end{center}

\end{document}